\chapter{交換対称性、スピン、そして二つの負符号}

\section{殻の充填とスピン--統計の問い}

殻模型の復習は、電子に付随する二つの性質から始まる。電子はスピン1/2をもち、その二つの無次元スピン射影は
\begin{equation}
m_s=+\frac12,
\qquad
m_s=-\frac12,
\end{equation}
であり、電子はパウリの排他原理に従う。これらの事実は関連しているが、同じ命題ではない。

``同じ状態''という表現には注意が必要である。二つの電子が同一の完全な一粒子状態を占めることはできない。それでも、二つの電子はスピン状態が異なれば一つの空間軌道を共有しうる。ヘリウムに対する独立電子近似では、両方の電子が最低の空間軌道を占め、一方はスピン上向き、他方はスピン下向きである。より正確には、``第一が上向きで第二が下向き''と``第一が下向きで第二が上向き''という選択肢は、もつれ合った二電子スピン状態に属する。

これは古い殻模型の近似にすぎない。電子間の相互作用を無視しており、その単純な数え上げのために用いている。リチウムでは、最初の二つの電子が最低の空間軌道を逆向きのスピンで満たす。第三の電子はそれらの完全な状態のいずれにも入ることができず、より高い軌道を占めなければならない。

次の段階はおなじみの軌道角運動量の図式である。
\begin{equation}
\ell=0,\qquad \ell=1,\qquad \ell=2
\end{equation}
である区分は\(S\)、\(P\)、\(D\)状態と呼ばれる。\(\ell\)の固定された値に対して
\begin{equation}
2\ell+1
\end{equation}
個の軌道状態があるので、
\begin{equation}
\ell=0\Rightarrow 1,
\qquad
\ell=1\Rightarrow 3,
\qquad
\ell=2\Rightarrow 5.
\end{equation}
クーロン問題には追加の縮退がある。異なる\(\ell\)区分からの動径準位が同じエネルギーをもちうる。殻の充填の議論では、次の殻は一つの\(S\)軌道と三つの\(P\)軌道を含む。したがって
\begin{equation}
1+3=4
\end{equation}
個の空間状態をもつ。スピンが数え上げを倍にする：
\begin{equation}
2(1+3)=8.
\end{equation}

これらの八つの完全な一電子状態は、リチウムからネオンへの単純な進行で満たされる。閉殻の先では、ナトリウムが再び一つの外殻電子をもち、それゆえこの粗い描像ではリチウムに似る。化学は目的ではない。その要点は、半整数スピンと排他がともに現れることである。

光子が対照を提供する。光子は整数スピンをもち、排他規則に従わない。古典電磁波は同じモードを占める多数の光子に対応する。ボソンは共通の状態に蓄積しうるが、同種のフェルミオンは一つの完全な状態を二度以上占めることはできない。問いは、これらの二つの組み合わせが別様でありえたかどうかである。すなわち、排他を伴わない半整数スピン、あるいは排他を伴う整数スピンである。三次元空間ではスピン--統計の関係が否と言う。おなじみの主張は一次元および二次元空間では限定を要し、場の理論的結果の証明はここでは試みない。

\begin{classroomqa}
\lecturetimestamp{00:11:48}
\audiencequestion{一つの状態を占めることのできるボソン的量子の最大数はあるか？}
\lecturerresponse{いいえ。ボソンには占有数の制限がなく、同じ状態に蓄積する傾向がある。}
\end{classroomqa}

\section{多粒子波動関数}

この関連は二つの負符号を通して接近する。一つは交換から、一つは回転からである。交換から始め、一時的にスピンを抑える。

\(x_i\)が粒子\(i\)の位置を表すとする。添字は粒子のラベルであり、デカルト成分ではない。三つの空間座標は単一の記号\(x_i\)に圧縮される。\(n\)粒子に対する位置基底は
\begin{equation}
\lvert x_1,x_2,\ldots,x_n\rangle.
\end{equation}
である。抽象状態\(\lvert\Psi\rangle\)に対して、座標波動関数は
\begin{equation}
\psi(x_1,x_2,\ldots,x_n)
=
\langle x_1,x_2,\ldots,x_n\mid\Psi\rangle.
\end{equation}
である。その絶対値の二乗、
\begin{equation}
\bigl\lvert\psi(x_1,x_2,\ldots,x_n)\bigr\rvert^2,
\end{equation}
は、粒子\(1\)を\(x_1\)に、粒子\(2\)を\(x_2\)に、以下同様に見出す同時確率密度である。

ここで粒子が同種であると仮定する。まず一方の電子を青、他方を赤に塗ったと想像する。すると``青が\(x_1\)に、赤が\(x_2\)''と、``青が\(x_2\)に、赤が\(x_1\)''を区別できたであろう。実際の電子はそのような色や個別のタグをもたない。帳簿上のラベルを入れ替えても新しい物理的配置は生じない。

全体位相は確率や期待値を変えない。確定した交換区分にある状態に対して、座標波動関数はしたがって
\begin{equation}
\psi(x_2,x_1)=\eta\,\psi(x_1,x_2),
\qquad
\lvert\eta\rvert=1.
\end{equation}
に従いうる。同じ交換を二度適用すると元の引数が戻り、
\begin{equation}
\eta^2=1,
\end{equation}
が得られるので、可能な二つの交換固有値は
\begin{equation}
\eta=+1
\qquad\hbox{or}\qquad
\eta=-1.
\end{equation}

\editorialnote{二重交換代数が可能な固有値を決定するのは、状態が確定した交換挙動を持つとされた後だけである。固定された種の許される状態が一つの交換セクターに属するというさらなる制限は物理的規則であり、\(\eta^2=1\)だけから従うものではない。同じ限定が以下の排他論証にも当てはまる。反対称性はフェルミオンセクターが選択された後に排他を意味するが、交換の代数それ自体はそのセクターを選択しない。}

ボソンに対して許される波動関数は対称であり、
\begin{equation}
\psi_B(x_1,x_2)=\psi_B(x_2,x_1),
\end{equation}
一方フェルミオンに対しては反対称である。
\begin{equation}
\psi_F(x_1,x_2)=-\psi_F(x_2,x_1).
\end{equation}
負号は状態全体に掛かる場合には見えないが、振幅が加えられる場合には重要である。こうして交換対称性は観測可能な帰結を得る。

\section{積、和、排他}

まず二つの粒子が同じ一粒子軌道\(\psi_0\)に割り当てられているとしよう。それらの積は
\begin{equation}
\psi_0(x_1)\psi_0(x_2).
\end{equation}
である。粒子座標を交換すると
\begin{equation}
\psi_0(x_2)\psi_0(x_1)
=
\psi_0(x_1)\psi_0(x_2).
\end{equation}
となる。波動関数の値は数であり演算子ではないので、因子は可換である。同一状態の積は対称である。それは許されるボソン形であるが、二つの同一フェルミオンを記述することはできない。

次に粒子を異なる一粒子軌道に割り当てる。
\begin{equation}
\psi_0(x_1)\psi_1(x_2).
\end{equation}
交換すると
\begin{equation}
\psi_0(x_2)\psi_1(x_1).
\end{equation}
が生じる。異なる関数に対しては、交換された積は一般に元の表式でもその負でもない。したがって裸の積はどちらの要求される交換対称性も持たない。これは一方の粒子が\(\psi_0\)を占め、他方が\(\psi_1\)を占めることを禁じるものではない。二粒子状態は二つの割り当てを組み合わせなければならないということを意味する。

対称結合は
\begin{equation}
\psi_B^{(01)}(x_1,x_2)
\propto
\psi_0(x_1)\psi_1(x_2)
+
\psi_0(x_2)\psi_1(x_1),
\end{equation}
であり、反対称結合は
\begin{equation}
\psi_F^{(01)}(x_1,x_2)
\propto
\psi_0(x_1)\psi_1(x_2)
-
\psi_0(x_2)\psi_1(x_1).
\end{equation}
である。交換論証には規格化が必要ないので、比例の符号で十分である。どちらの状態も各軌道に一つの粒子を含む。それらは二つの区別できない割り当てがどのように組み合わされるかだけが異なる。

ここで軌道を一致させよう。ボソン結合は
\begin{align}
\psi_B^{(00)}(x_1,x_2)
&\propto
\psi_0(x_1)\psi_0(x_2)
+
\psi_0(x_2)\psi_0(x_1)\\
&=
2\psi_0(x_1)\psi_0(x_2),
\end{align}
となり、これは非零である。フェルミオン結合は
\begin{align}
\psi_F^{(00)}(x_1,x_2)
&\propto
\psi_0(x_1)\psi_0(x_2)
-
\psi_0(x_2)\psi_0(x_1)\\
&=0.
\end{align}
を与える。反対称交換規則のもとで排他が従う。すなわち、両方のフェルミオンが同じ一粒子状態にある非零の反対称状態は存在しない。

多粒子状態に対して、この規則は二つの引数の任意の置換に適用される。置換はボソン波動関数を変えず、フェルミオン波動関数の符号を反転させる。同じ置換を繰り返すことはどちらの場合も自明である。驚くべき点は、一回の交換が観測可能な粒子ラベルを変えないにもかかわらず、フェルミオン状態の符号を変えることである。

\section{完全な粒子ラベルにおけるスピン}

ここでスピンを小さな変更とともに復元できる。一粒子振幅は位置と離散的なスピンラベルとの両方に依存し、
\begin{equation}
\psi(x,\sigma),
\end{equation}
ここで\(\sigma=+1\)または\(-1\)が選ばれた\(\sigma_z\)の結果をラベルする。その絶対値の二乗は、述べられた位置とスピンラベルに対する結合確率密度である。

多粒子に対して波動関数は
\begin{equation}
\psi(x_1,\sigma_1;\,x_2,\sigma_2;\,\ldots).
\end{equation}
である。交換は二つの粒子に属するすべてを入れ替え、その位置だけではない。
\begin{align}
\psi_B(x_1,\sigma_1;\,x_2,\sigma_2)
&=
+\psi_B(x_2,\sigma_2;\,x_1,\sigma_1),\\
\psi_F(x_1,\sigma_1;\,x_2,\sigma_2)
&=
-\psi_F(x_2,\sigma_2;\,x_1,\sigma_1).
\end{align}
これが冒頭の殻模型の記述の正確な形である。二つの電子が一つの空間軌道を共有できるのは、スピンを含むそれらの完全な状態が異なる場合だけである。

\begin{classroomqa}
\lecturetimestamp{00:34:33}
\audiencequestion{交換対称性は置換群の一例か。}
\lecturerresponse{はい。粒子を交換することは置換であり、プラスまたはマイナスの変換則は置換群の表現を与える。}
\end{classroomqa}

\begin{classroomqa}
\lecturetimestamp{00:35:12}
\audiencequestion{なぜ二粒子の構成は一粒子波動関数の和ではなく積なのか。}
\lecturerresponse{和は一つの粒子が重ね合わせにあることを記述する。積は二粒子テンソル積空間に属し、各一粒子状態に一つの粒子が割り当てられていることを記述する。}
\end{classroomqa}

完全ラベルの交換を下の板書に示す。\footnote{位置--スピンラベルの順序は、判読可能な黒板のフレームと00からの周囲の説明から再構成される。33：43から00：34：25。}
\begin{figure}[htbp]
\centering
\includegraphics[width=0.9\linewidth]{lecture_05_frame_03.png}
\caption{二つの同一粒子を交換すると、それらの完全な位置--スピンラベルが入れ替わる。ボソン状態は変わらず、フェルミオン状態は符号を変える。}
\lectureframe{00:35:25}
\end{figure}

これらは線形性の異なる用法である。一粒子に対して、
\begin{equation}
\alpha\psi_0(x)+\beta\psi_1(x)
\end{equation}
は二つの選択肢の重ね合わせである。二粒子に対して、
\begin{equation}
\psi_0(x_1)\psi_1(x_2)
\end{equation}
は二つの粒子座標を持つ空間上の積振幅である。対称化または反対称化はその積を交換された相手と組み合わせる。

\begin{classroomqa}
\lecturetimestamp{00:36:13}
\audiencequestion{粒子を交換すると干渉効果が生じるのか。}
\lecturerresponse{はい。同一状態のフェルミオン的振幅が消えることは、交換された二つの寄与の間の破壊的干渉として理解できる。}
\end{classroomqa}

\begin{classroomqa}
\lecturetimestamp{00:37:13}
\audiencequestion{パウリはこの同じ対称性の議論を用いて排他律を導いたのか。}
\lecturerresponse{パウリは経験的事実から出発し、それを表現できる数学的枠組みを見出した。彼は当初、より基礎的な理論から排他律を導いたのではない。}
\end{classroomqa}

反対称性は、経験的な排他律が表現される数学的枠組みを与える。

\begin{classroomqa}
\lecturetimestamp{00:39:01}
\audiencequestion{粒子交換の干渉は、代替経路間の干渉と関係しているのか。}
\lecturerresponse{はい。一方の履歴では二つの粒子は交差せずに通過し、他方の履歴ではそれらは交差して交換された行き先へと向かう。振幅はボソンでは加算され、フェルミオンでは相殺しうる。}
\end{classroomqa}

交換が第一の負符号を与える。回転が第二の負符号を与える。

\section{第二のマイナス符号：\texorpdfstring{\(2\pi\)}{2 パイ}}

完全な回転は何もしていないように見える。あらゆる通常の幾何学的方向は角度 \(2\pi\) の後に出発点へ戻る。しかし戻らなければならないのは確率と期待値である。状態ベクトル自体は、交換の下でそうでありうるように、全体的な位相を伴って戻ってよい。

\begin{classroomqa}
\lecturetimestamp{00:42:12}
\audiencequestion{回転にはどの軸が用いられ、生成子は軌道角運動量、スピン、それとも両方なのか。}
\lecturerresponse{外部軸を選んでよい。以下ではそれを \(z\) 軸とする。生成子は軌道寄与とスピン寄与の両方を含む全角運動量である。スピン自体は選ばれた軸に沿って向く必要はない。}
\end{classroomqa}

\editorialnote{この議論における角運動量の公式は \(\hbar=1\) を用いるので、\(m\) は \(J_z\) の無次元固有値である。}

\(\psi(\theta)\) を、回転パラメータが角度 \(\theta\) を通って進むにつれて得られる状態を表すものとする。ここで用いる規約では、
\begin{equation}
J_z\psi(\theta)
=
-i\frac{\partial}{\partial\theta}\psi(\theta).
\end{equation}
確定的な \(J_z\) をもつ状態に対して、
\begin{equation}
J_z\psi(\theta)=m\psi(\theta).
\end{equation}
これらの方程式を組み合わせると、
\begin{equation}
-i\frac{\partial}{\partial\theta}\psi(\theta)
=
m\psi(\theta),
\end{equation}
となり、解は
\begin{equation}
\psi(\theta)=e^{im\theta}\psi(0).
\end{equation}

\(m\) の許される値は整数または半整数である。したがって
\begin{equation}
\psi(2\pi)
=
\begin{cases}
+\psi(0), & m\in\mathbb Z,\\[4pt]
-\psi(0), & m\in\mathbb Z+\tfrac12.
\end{cases}
\end{equation}
さらにもう一回転の後、
\begin{equation}
\psi(4\pi)=\psi(0)
\end{equation}
いずれの場合でも成り立つ。これが第二の負符号である。半整数角運動量多重項の状態は \(2\pi\) 回転の下で符号を変える。

この符号は \(\lvert\psi\rvert^2\) を変えないので、状態全体に掛かる場合には見えない。それでもなお、\(2\pi\) 回転は半整数スピンの状態上で恒等変換によって表されるのではない。この区別は、符号が別のコヒーレントな振幅に対して相対的にされる場合に観測可能となる。

\section{\texorpdfstring{\(2\pi\) と \(4\pi\)}{2 パイおよび 4 パイ} の経路}

\(2\pi\) と \(4\pi\) の区別は、これ以上の代数なしに可視化できる。片手でカップを持ち、それを放さずに一回転させる。カップは元の向きに戻るが、腕はねじれている。さらに二回目の回転を続けると、カップの運動を逆転させることなく腕を元の位置に戻すことができる。

関連する構成では、球を箱の中に置き、ひもで箱に結びつける。球を \(2\pi\) 回転させるともつれが残る。通常の三次元空間では、\(4\pi\) 回転に伴うもつれはほどくことができる。問題となる対象は、球の最終的な向きだけでなく、回転の空間を通ってたどられた経路である。

\begin{classroomqa}
\lecturetimestamp{00:54:22}
\audiencequestion{球を同じ軸のまわりで二回転させると、ひもは本当にほどけるのか。それには別の空間次元が必要なのか。}
\lecturerresponse{\(4\pi\) の経路は通常の三次元空間でほどくことができる。時間は運動をパラメータ化するのであり、別の空間方向ではない。\(2\pi\) の経路はこのようにはほどけず、\(4\pi\) の経路はほどける。}
\end{classroomqa}

これらの実演は、回転群の形式的理論を展開することなく、非自明な \(2\pi\) 経路とほどける \(4\pi\) 経路を区別する。

\section{ベルト描像における交換と回転}

二つの負符号が今や並んでいる。二つの同一なフェルミオンを交換すると、それらの結合波動関数の符号が変わる。一つの半整数スピン状態を \(2\pi\) 回転させても、その符号が変わる。ベルトは二つの操作の間の幾何学的関係を与える。

ベルトの下端二つが交換されるようにベルトを配置する。別の見方をすれば、同じベルトは\(2\pi\)のねじれを持つ。交換を\(2\pi\)の回転と一緒に行うとねじれが取り除かれ、ねじれのないベルトが残る。この結合過程は連続的に何もしないことに還元できる。\footnote{交換--回転の観察は講義ではDavid Finkelsteinに帰されている。原記録には特定の出版物が同定されていないため、ここではその帰属としてのみ名前を保持する。}

\begin{classroomqa}
\lecturetimestamp{01:01:12}
\audiencequestion{これは粒子のトポロジーについての主張ではなく、ベルトが空間にどのように埋め込まれているかについての主張にすぎないのか。}
\lecturerresponse{ベルトは時空における粒子の履歴として読むことができる。時間を上向きに取ると、その断面は粒子のような対象を与える。対が生成され、一方の対象が\(2\pi\)回転され、二つの同一な対象が交換される。その履歴では回転と交換の位相が相殺する。}
\end{classroomqa}

この履歴は位相の帳簿付けを明示する。粒子が無い状態から始める。二つの離れた場所で粒子--反粒子対を生成し、中間の時刻に四つの粒子を残す。一方の粒子を\(2\pi\)回転させながら、配置の別の場所で二つの同一な粒子を交換する。完全な履歴は、対が生成され後に消滅する初等的な履歴に変形できる。ここで議論している符号の帳簿付けでは、その自明な履歴には追加の位相がない。\(\eta_{\mathrm{rot}}\)と\(\eta_{\mathrm{ex}}\)が二つの符号を表すなら、口頭で述べられた相殺は簡潔に次のように書ける\footnote{この位相方程式は01から述べられた相殺を記録している:06:50から01:07:18; 講義はこの構成を厳密な証明というよりも示唆的なものとして提示している。}
\begin{equation}
\eta_{\mathrm{rot}}\eta_{\mathrm{ex}}=1.
\end{equation}

\begin{classroomqa}
\lecturetimestamp{01:03:34}
\audiencequestion{二つの粒子のそれぞれを\(\pi\)動かし、その後それらを再配向させることで、交換を回転に関係づけることができるか。}
\lecturerresponse{その構成は関連する数学を反映しているかもしれないが、二つの別々の\(\pi\)回転を用いている。ここで分析している過程は一つの\(2\pi\)回転を含むので、二つの議論は同一ではない。}
\end{classroomqa}

この議論は粒子のような対象を生成し消滅させる可能性に依存している。粒子が永遠であるような理論ではこの形をとらないであろう。またこれはスピン--統計定理の標準的な証明として提示されているものでもない。その目的は、交換と回転の符号が整合するトポロジー的な理由を示すことにある。

\section{塊、鏡像、および消滅}

この構成は局在した場の配置について繰り返すことができる。一つまたは複数の場がほとんど至る所で真空値をとるが、小さな領域で非零の形と構造をもつとしよう。その局在した配置が塊である。

次に、第一の適切な鏡像として第二の配置を作る。一方を粒子、他方を反粒子と呼ぶ。それらを近づける。場の配置の対応する部分が出会い、互いを打ち消す。中間の段階では、塊の一部はすでに消えているが他の部分は残っているかもしれない。変形を続けると残りの構造が取り除かれ、場は真空配置に戻る。ベルトの二つの逆向きのねじれも同じ描像を与える: ねじれを押し合わせるとそれらはほどける。

この過程は逆向きに走らせることができる。空の空間は塊とその鏡像の相棒を生成しうる。そのような対を二つ生成し、一方の塊を\(2\pi\)回転させ、二つの同一な塊を交換し、再び対を消滅させてもよい。そのとき同じ交換--回転の相殺が適用される。これは構成のヒューリスティックな拡張であり続け、あらゆる局在励起についての定理ではない。

\begin{classroomqa}
\lecturetimestamp{01:11:30}
\audiencequestion{これらの塊のような場の配置には物理的な実現があるか。}
\lecturerresponse{はい。例としてはSkyrmion型の記述における原子核、磁気単極子、その他の場の粒子のような励起が含まれる。}
\end{classroomqa}

鏡像の塊の描像が一般化を理解可能にするものである。関連する特徴は、対象が初めに素な点粒子であったということではなく、粒子のような励起とその相棒が連続的に生成され、近づけられ、消滅されうるということである。その特徴が回転的振る舞いと交換的振る舞いの間のより広い関係を動機づける。

\begin{classroomqa}
\lecturetimestamp{01:16:12}
\audiencequestion{スピン1/2とは、状態が通常の回転位相の半分しか蓄積しないという意味でしょうか。}
\lecturerresponse{この規約での位相は\(e^{im\theta}\)です。半整数の\(m\)に対して、\(2\pi\)回転は\(-1\)を与え、一方\(4\pi\)回転は\(+1\)を与えます。}
\end{classroomqa}

\section{大域的符号とヌル実験}

\(2\pi\)回転によるマイナス符号は観測できるか。まず、それが観測できない実験から始める。電子を磁場中に閉じ込め、装置をゆっくりと\(2\pi\)回転させながらそのスピンを磁場に追随させ、その後電子を干渉装置へ放出する。これを、回転を行わない同じ手順と比較する。

回転は状態を
\begin{equation}
\psi\longmapsto-\psi,
\end{equation}
だけ変化させるが、マイナス符号は状態全体にかかる。すべての確率は変わらないので、最終的な干渉縞には予備的な回転が行われたかどうかの記録は残らない。これはヌル実験である。共通位相には干渉すべき相手がない。

\begin{classroomqa}
\lecturetimestamp{01:20:01}
\audiencequestion{なぜ磁場の回転はゆっくり行わなければならないのですか。}
\lecturerresponse{スピンは回転する磁場に追随し続けなければならない。磁場をあまりに速く動かすと、スピンは磁場に束縛されたままでいられない。}
\end{classroomqa}

\section{二つのコヒーレントな枝の中の一つの電子}

符号は、それを相対的なものにしたときに観測可能になる。一つの電子をビームスプリッターに通す。その波動関数の一部は一方の経路を進み、一部は他方の経路を進む。
\begin{equation}
\psi(x)=\psi_1(x)+\psi_2(x).
\end{equation}
これは一つの電子であり、二つの電子ではない。二つの選択肢のコヒーレントな重ね合わせである。この区別は明確に保たなければならない。二つの波束は二つの独立な粒子ではない。

二つの波束を、磁場を含む別々の領域へ導く。一方の領域で磁場を\(2\pi\)回転させ、他方は変えない。半整数スピンの状態に対して、
\begin{equation}
\psi_1(x)+\psi_2(x)
\longmapsto
\psi_1(x)-\psi_2(x),
\end{equation}
ここで回転した枝のラベルは\(2\)と選ばれている。符号は今や相対的である。二つの波束が後に重なるとき、加算は減算に変わっている。建設的干渉の点が破壊的干渉の点に、あるいは節にさえなりうる。

\begin{classroomqa}
\lecturetimestamp{01:23:46}
\audiencequestion{電子が一つしかなく、波動関数の枝が二つあるとき、何が回転しているのですか。}
\lecturerresponse{電子は二つの局在した選択肢のコヒーレントな重ね合わせの中にある。一方の枝は回転され、他方は回転されない。どちらの選択肢が起こったかを知るために箱を開けると、その測定は干渉を破壊する。}
\end{classroomqa}

箱を開ければ、一方の箱か他方の箱にある一つの電子だけが明らかになるだろう。その測定がなければ、両方の振幅が状態に残る。実験を繰り返すと、干渉縞が一回の検出ごとに蓄積されていく。相対的なマイナス符号によって生じた節では、電子は検出されない。

\section{磁気歳差と二重スリット干渉}

より実用的な概略的配置は、非常に遅い箱の実験を置き換える。コヒーレントな電子波を二つの経路に分ける。一方の経路では、電子を、そのスピンが初め磁場に垂直であるような磁気領域を通らせる。するとスピンは歳差運動する。磁場の強さと通過時間を選ぶことにより、その枝のスピンは一回の完全な\(2\pi\)回転を行い、他方の枝はそのような回転を全く受けない。\footnote{講義は、関連するスピノル位相の実験が行われたと述べ、その原理を説明しているが、資料記録は論文や特定の装置を同定していない。したがってここでの記述は、講義で与えられた構成に限られる。}

\begin{classroomqa}
\lecturetimestamp{01:27:05}
\audiencequestion{単一電子の経路はどのようにして決定できるのですか。}
\lecturerresponse{実験を台無しにすることなく決定することはできない。電子がどちらの経路を通ったかを記録するいかなる測定も、干渉を破壊する。}
\end{classroomqa}

\begin{classroomqa}
\lecturetimestamp{01:27:50}
\audiencequestion{波動関数の枝とは正確には何であり、スプリッターはどのようにして二つの枝を生み出すのですか。}
\lecturerresponse{枝とは、一方の経路に沿って局在した波束である。スプリッターは入射波を空間的に分離した二つの波束に変え、電場または磁場がそれらの波束を異なる経路に沿って導くことができる。}
\end{classroomqa}

磁気領域を除けば、これは二重スリット実験である。入射電子波は二つの波束に分かれる。出ていく波は共通の検出領域に広がり、そこでそれらの振幅が重なる。相対的回転がなければ、中心の振幅は加算されうる。一方の経路での\(2\pi\)スピン回転は一方の振幅を反転させ、同じ寄与を相殺させうる。

\begin{classroomqa}
\lecturetimestamp{01:30:46}
\audiencequestion{電子が別々の管やスリットを通るのなら、なぜ干渉が起こりうるのでしょうか。}
\lecturerresponse{最終的な検出は経路を明らかにしません。二つの出ていく波の振幅は広がり、検出器で重なり合います。どちらの経路を通ったかを見ようとすれば、干渉は破壊されるでしょう。}
\end{classroomqa}

\begin{classroomqa}
\lecturetimestamp{01:33:22}
\audiencequestion{管やスリットの直径は電子の波長と同程度でなければならないのでしょうか。}
\lecturerresponse{文字通り回折に基づく二重スリットの実現では、必要な広がりを生じさせるために開口部は十分に小さくなければなりません。ここで述べられている装置は実際的な設計というよりも模式的なものです。}
\end{classroomqa}

\begin{classroomqa}
\lecturetimestamp{01:33:44}
\audiencequestion{光子について対応する\(2\pi\)の偏光回転を行えば、干渉縞は変わらないのでしょうか。}
\lecturerresponse{はい。光子は整数スピンをもつので、完全な\(2\pi\)回転はフェルミオン的な負符号をもたらしません。}
\end{classroomqa}

\begin{classroomqa}
\lecturetimestamp{01:34:12}
\audiencequestion{実際の実験が理想化された版からわずかに逸脱する点は何でしょうか。}
\lecturerresponse{理想的な構成ではスピンを断熱的に回転させます。実際的な記述では、電子は磁場領域を通り抜けてそのスピンはより速く歳差運動しますが、測定される位相は依然として\(2\pi\)回転に対応する位相です。}
\end{classroomqa}

遅い構成は、スピンが余分な運動を伴わずに場に従うことを保証することで、この考えを切り出します。歳差運動の構成は、同じ相対符号に動的に到達します。どちらの版でも、干渉縞が現れる前に、二つの経路は区別できなくなっていなければなりません。

\begin{classroomqa}
\lecturetimestamp{01:36:15}
\audiencequestion{二つの枝はスプリッターの後で再結合されなければならないのでしょうか。}
\lecturerresponse{はい。出てくる波束は重なり合うか、共通の検出領域に広がって、その振幅が干渉縞を形成できるようにしなければなりません。}
\end{classroomqa}

\begin{classroomqa}
\lecturetimestamp{01:38:12}
\audiencequestion{この一電子回転実験はパウリ排他律の帰結なのでしょうか。}
\lecturerresponse{いいえ。電子は一つだけです。この実験は回転の下でのスピノル位相を測定するものであり、二つのフェルミオンが一つの状態から排除されることを測定するものではありません。}
\end{classroomqa}

\section{交換干渉と相対回転}

交換符号は同じ一般的な方法で露わにできるかもしれませんが、提案された装置はより複雑です。二つの電子から始めて、それらの結合波動関数を二つのコヒーレントな枝に分けます。依然として電子は二つだけで、四つではありません。対は一方の選択肢か他方に従います。一方の枝では二つの電子を交換し、他方の枝では交換しないままにします。枝を再び合わせた後、フェルミオン的交換符号は相対位相として現れ、干渉を変えうるのです。

最終的な提案は二つの負符号を直接比較します。一方の枝はフェルミオン的交換位相を獲得し、他方の枝は対応する\(2\pi\)回転位相を獲得します。どちらの位相も\(-1\)であるため、相対的な効果は相殺されます。\footnote{01:40:26と01:41:15の間の記録は、交換位相、\(2\pi\)回転位相、およびそれらの相殺を明確に与えていますが、回転の詳細な割り当ては利用可能な発話からは復元できません。} 講義はこれを提案された思考実験として提示し、講演者の知る限りそれは実行されていないと述べています。

\begin{classroomqa}
\lecturetimestamp{01:41:34}
\audiencequestion{交換位相を観測することは、二つの電子の区別不可能性に挑戦することになるのでしょうか。}
\lecturerresponse{いいえ。量子的な区別不可能性は古典的な同一性よりも微妙です。それは、いかなる観測可能なラベルも粒子を区別しないにもかかわらず、状態ベクトルが交換の下で符号を変えることを許します。}
\end{classroomqa}

最後に一つの限定があります。孤立した世界全体を\(2\pi\)回転させることは共通の位相だけを与え、いかなる内部実験もそれを明らかにできません。検出可能な実験は、一つのコヒーレントな枝を別の枝に対して回転させます。したがって観測されるのは絶対的な符号ではなく、相対的な符号によって生じる干渉です。
